\section{Method}
\label{sec:method}

We introduce the \textsc{ComposeHead}, a trainable module that represents candidate
solutions as sums of primitive compositions, together with a search-and-refine
pipeline that converts learned parameters into exact symbolic expressions.
The key design principle is to embed the separable product structure common to
analytical PDE solutions directly into the architecture, so that the optimizer
need only recover scalar coefficients rather than discover structural patterns
from scratch.

%% ----------------------------------------------------------------
\subsection{ComposeHead Architecture}
\label{sec:composehead}

A \textsc{ComposeHead} with $n$ input dimensions represents a function
\begin{equation}
\label{eq:composehead}
    f(\mathbf{x}) \;=\; \beta \;+\; \sum_{i=1}^{T} c_i \, \phi_i(\mathbf{x}),
\end{equation}
where $\beta \in \mathbb{R}$ is a trainable bias, $c_i \in \mathbb{R}$ are trainable
coefficients, and each $\phi_i$ is a \emph{term}---a composition of primitives drawn
from a function basis
$\mathcal{F} = \{\sin, \cos, \exp, \ln, \mathrm{id}, \mathrm{sq}, \mathrm{inv}\}$.
All primitives are implemented via verified EML constructions (Section~\ref{sec:background}),
ensuring that the entire forward pass traces back to the base operation
$\mathrm{eml}(x, y) = e^x - \ln y$.

We define three term types of increasing structural commitment.

\paragraph{PrimitiveTerm.}
A single primitive applied to a learned linear combination of inputs:
\begin{equation}
\label{eq:primitiveterm}
    \phi(\mathbf{x}) = g\!\left(\mathbf{w}^\top \mathbf{x} + b\right),
    \quad g \in \mathcal{F},
\end{equation}
with trainable parameters $\mathbf{w} \in \mathbb{R}^n$, $b \in \mathbb{R}$, and
coefficient $c \in \mathbb{R}$.  The weight vector $\mathbf{w}$ allows the term to
attend to arbitrary linear combinations of input dimensions, providing maximum
flexibility at the cost of a larger parameter space.

\paragraph{ProductTerm.}
A pairwise product of two PrimitiveTerms sharing a single coefficient:
\begin{equation}
\label{eq:productterm}
    \phi(\mathbf{x}) = g_1\!\left(\mathbf{w}_1^\top \mathbf{x} + b_1\right)
    \cdot g_2\!\left(\mathbf{w}_2^\top \mathbf{x} + b_2\right),
    \quad g_1, g_2 \in \mathcal{F}.
\end{equation}
Each factor maintains its own weight vector, enabling expressions such as
$\sin(\mathbf{w}_1^\top \mathbf{x}) \cos(\mathbf{w}_2^\top \mathbf{x})$.
In practice, the optimizer must learn near--one-hot patterns in $\mathbf{w}_1$ and
$\mathbf{w}_2$ to isolate individual input dimensions---a requirement that
introduces spurious local minima when $n$ is even moderately large.

\paragraph{SeparableTerm (key innovation).}
A product of $K$ axis-aligned factors, each operating on exactly one input dimension:
\begin{equation}
\label{eq:separableterm}
    \phi(\mathbf{x}) = \prod_{k=1}^{K}
    g_k\!\left(a_k \, x_{d_k} + b_k\right),
    \quad g_k \in \mathcal{F},\;\;
    d_k \in \{1, \ldots, n\},
\end{equation}
where each \emph{AxisTerm} factor has only two trainable parameters: a scale
$a_k \in \mathbb{R}$ and a bias $b_k \in \mathbb{R}$.  The entire SeparableTerm
shares a single coefficient $c \in \mathbb{R}$.

\smallskip

The central observation motivating the SeparableTerm is that the vast majority of
known analytical PDE solutions are \emph{separable}---expressible as products of
univariate functions of individual coordinates.  Classical examples include
normal modes of the wave equation, Fourier--series solutions to the heat equation,
and the Taylor--Green vortex solution to the Navier--Stokes equations.
By making this product-of-univariates structure explicit in the architecture,
the optimizer is relieved of the burden of discovering one-hot attention patterns
in free weight vectors.  It need only recover the scalar scales and biases
within each factor, a dramatically simpler optimization landscape.

Table~\ref{tab:param_counts} summarizes the parameter counts.  The SeparableTerm
is notable in that its parameter count depends only on the number of factors $K$
and is \emph{independent} of the input dimensionality $n$.

\begin{table}[t]
\centering
\caption{Parameter counts per term type for $n$ input dimensions.  The SeparableTerm
with $K$ factors is independent of $n$, yielding a favorable scaling for
high-dimensional separable targets.}
\label{tab:param_counts}
\smallskip
\begin{tabular}{lcc}
\toprule
\textbf{Term type} & \textbf{Parameters} & \textbf{Scaling} \\
\midrule
PrimitiveTerm     & $n + 2$   & $\mathcal{O}(n)$ \\
ProductTerm       & $2n + 3$  & $\mathcal{O}(n)$ \\
SeparableTerm ($K$ factors) & $2K + 1$  & $\mathcal{O}(K)$, independent of $n$ \\
\bottomrule
\end{tabular}
\end{table}

%% ----------------------------------------------------------------
\subsection{Training Pipeline}
\label{sec:pipeline}

Given data $(\mathbf{X}, \mathbf{y})$ with $\mathbf{X} \in \mathbb{R}^{N \times n}$,
we seek a symbolic expression $\hat{f}$ that approximates the mapping
$\mathbf{x} \mapsto y$.  Algorithm~\ref{alg:discovery} presents the full
discovery pipeline.

\begin{algorithm}[t]
\caption{Compositional EML Discovery}
\label{alg:discovery}
\begin{algorithmic}[1]
\REQUIRE Data $(\mathbf{X}, \mathbf{y})$, input dimension $n$,
         function basis $\mathcal{F} = \{\sin, \cos, \exp, \ldots\}$
\ENSURE  Symbolic expression $\hat{f}$

\medskip
\STATE \textbf{Candidate enumeration.}
    For each assignment of $K$ functions from $\mathcal{F}$ to $K$ input
    dimensions $(d_1, \ldots, d_K)$, instantiate a SeparableTerm
    $\phi = \prod_{k=1}^{K} g_k(a_k \, x_{d_k} + b_k)$.
    \hfill \COMMENT{$|\mathcal{F}|^K$ candidates for fixed dims}

\medskip
\STATE \textbf{Parallel evaluation.}
    Train each candidate independently for $N_{\mathrm{search}}$ steps using
    the Adam optimizer~\citep{paszke2019}, minimizing
    $\mathcal{L} = \frac{1}{N}\sum_{i=1}^{N}
    \bigl(\hat{y}_i - y_i\bigr)^2$.
    Rank candidates by final MSE.

\medskip
\STATE \textbf{Selection.}
    Retain the top-$1$ candidate (or top-$M$ for multi-term targets).

\medskip
\STATE \textbf{Fine-tuning.}
    Train the selected term(s) for $N_{\mathrm{fine}}$ additional steps at
    a reduced learning rate.

\medskip
\STATE \textbf{Normalization.}
    Canonicalize the learned parameterization:
    \begin{itemize}
        \item \emph{Exponential factors:}
            absorb the bias into the coefficient via
            $c \cdot \exp(ax + b) \;\to\; (c \, e^{b}) \cdot \exp(ax)$,
            then set $b \leftarrow 0$.
        \item \emph{Trigonometric factors:}
            reduce the phase $b$ to $[-\pi, \pi]$;
            absorb sign flips using $\sin(x + \pi) = -\sin(x)$ and
            $\cos(x + \pi) = -\cos(x)$.
        \item \emph{Negative coefficients:}
            if $c < 0$ and a $\sin$ factor is present, exploit
            $\sin(-x) = -\sin(x)$ to make $c$ positive by negating
            the scale and bias of that factor.
    \end{itemize}

\medskip
\STATE \textbf{Snapping (optional).}
    Round each parameter to the nearest element of a target set
    $\mathcal{S} = \{0, \pm 1, \pm \tfrac{1}{2}, \pm \tfrac{1}{3},
    \pm \tfrac{\pi}{2}, \pm \pi, \pm e, \pm \sqrt{2}, \pm \ln 2, \ldots\}$
    within tolerance $\tau$.
    If validation data are available, revert any snap that causes the loss
    to exceed a maximum ratio $r_{\max}$ relative to the pre-snap loss.

\medskip
\RETURN Symbolic expression extracted via SymPy conversion of the
    snapped parameters.
\end{algorithmic}
\end{algorithm}

For concreteness, consider the Taylor--Green vortex with inputs
$\mathbf{x} = (x, y, t)$ and basis $\mathcal{F} = \{\sin, \cos, \exp\}$.
The candidate enumeration in Step~1 generates $|\mathcal{F}|^3 = 27$ triple
SeparableTerms (one per function assignment to the three input dimensions).
Each candidate is trained for $N_{\mathrm{search}} = 2{,}000$ steps (Step~2),
the best is selected (Step~3), and fine-tuned for $N_{\mathrm{fine}} = 5{,}000$
steps (Step~4).  The entire search completes in under two minutes on a single
CPU, as each candidate contains only $2K + 1 = 7$ trainable parameters.

The normalization step (Step~5) is critical for reliable snapping.
The parameterization $c \cdot g(ax + b)$ admits gauge freedoms:
for exponential factors, the coefficient $c$ and the term $e^b$ can trade off
arbitrarily without changing the function value, creating a ridge in the loss
landscape along which the optimizer may drift.  Normalization breaks these
symmetries, concentrating the ``meaning'' of each parameter and ensuring that
the subsequent snapping step compares like with like.

%% ----------------------------------------------------------------
\subsection{Why This Works}
\label{sec:why}

We identify five properties of the proposed approach that collectively
explain its effectiveness on the class of problems we consider.

\paragraph{Separable ansatz matches target structure.}
The solutions to the PDEs we target---heat equation, wave equation,
Navier--Stokes in the Taylor--Green regime---are products of univariate functions
of the spatial and temporal coordinates.  The SeparableTerm directly
encodes this product-of-univariates structure, so that a correct solution
lies at a point in parameter space rather than in a complex submanifold
of a higher-dimensional space.

\paragraph{Reduced search space.}
Classical symbolic regression methods such as genetic programming~\citep{cranmer2023}
search over arbitrary expression trees, yielding a combinatorially large
search space.  Our approach instead enumerates function assignments to input
dimensions, producing $|\mathcal{F}|^K$ candidates---polynomial in
$|\mathcal{F}|$ for fixed $K$, and polynomial in $K$ for fixed $|\mathcal{F}|$.
For the problems considered here, $|\mathcal{F}| \leq 7$ and $K \leq 3$,
so exhaustive enumeration is tractable.

\paragraph{Normalization resolves gauge freedom.}
The parameterization $c \cdot g(ax + b)$ introduces redundancies.
For $g = \exp$, the identity $c \cdot \exp(ax + b) = (c\,e^{b}) \cdot \exp(ax)$
means that $c$ and $b$ are not independently identifiable.
For $g = \sin$ or $g = \cos$, the periodicity and parity identities
create analogous degeneracies.  The normalization step
(Algorithm~\ref{alg:discovery}, Step~5) systematically eliminates
these gauge freedoms, yielding a canonical form in which each parameter
has a unique interpretation.  This is essential: without normalization,
a parameter with true value $a = 1$ might be learned as $a = 0.93$
with compensating $b$ and $c$ values, causing the snapping step to fail.

\paragraph{Single-term search avoids signal splitting.}
When a multi-term model is trained end-to-end, the optimizer may distribute
the target signal across many terms, each capturing a small fraction
of the variance.  The resulting coefficient magnitudes are small and
difficult to distinguish from noise, making pruning unreliable.
By training each candidate SeparableTerm in isolation (Step~2 of
Algorithm~\ref{alg:discovery}), we force the optimizer to concentrate
all model capacity in a single term.  The correct candidate achieves
near-zero loss; incorrect candidates plateau at high loss.
This binary signal---correct versus incorrect---is far easier to threshold
than the continuous coefficient magnitudes in a multi-term model.

\paragraph{Gradient-based optimization.}
Unlike genetic programming approaches, every operation in our pipeline
is differentiable.  Each EML primitive---$\exp$ via $\mathrm{eml}(x, 1)$,
$\sin$ and $\cos$ via the imaginary and real parts of
$\mathrm{eml}(ix, 1) = e^{ix}$---supports standard backpropagation.
This allows us to leverage mature gradient-based optimizers (Adam) with
well-understood convergence properties, rather than relying on stochastic
search heuristics.  The per-candidate optimization problem is low-dimensional
($2K + 1$ parameters), smooth, and---for the correct function assignment---typically
convex in the neighborhood of the true solution.
